\documentclass[journal,twoside]{IEEEtran}
\usepackage[utf8]{inputenc}
\usepackage{cite}
\usepackage{amsmath,amssymb,amsfonts}
\usepackage{algorithmic}
\usepackage{graphicx}
\usepackage{textcomp}
\usepackage{url}
\usepackage{hyperref}
\usepackage{array}
\usepackage{booktabs}
\usepackage{amsmath}

\begin{document}

\title{A Physically-Constrained, Multi-Sensor, Temporally-Aware Deep Learning System for Recall-Optimized Multi-Hazard Prediction in Nepal at 500m Resolution}

\author{Kabish Bhattarai, \textit{BSc CSIT}, Tribhuvan University, Nepal}

\markboth{Technical Research Report, April 2026}%
{Bhattarai \MakeLowercase{\textit{et al.}}: Nepal Multi-Hazard prediction}

\maketitle

\begin{abstract}
Nepal's diverse topography and monsoon-dominated climate create a highly volatile environment susceptible to landslides and flash floods. This research presents a high-resolution (500m), physically-constrained deep learning framework for environmental hazard prediction. We transition from coarse regional averages to a localized point-based sampling architecture, utilizing 873 official disaster records from the Nepal BIPAD portal as ground-truth validation. The system fuses multi-sensor satellite telemetry (Sentinel-1, 2, 5P, and MODIS) with static topographic covariates including Slope, Aspect, and Topographic Position Index (TPI) derived from 30m SRTM data. A dual-strategy multi-scale spatial harmonization scheme prevents aliasing artifacts by spatially averaging high-resolution imagery over 500m footprints while bilinear-interpolating coarse climate fields upward. A Bi-directional Long Short-Term Memory (Bi-LSTM) network is trained on 7-day multi-channel sequences to capture antecedent moisture triggers, producing a continuous Environmental Hazard Risk Score $\hat{P}(H | \mathbf{X}) \in [0,1]$. Comparative evaluation against Random Forest, Gradient Boosting, and unidirectional LSTM baselines demonstrates that the Bi-LSTM achieves the highest hazard recall (0.81) with a ROC-AUC of 0.933, making it operationally preferable for early warning systems where minimizing missed disasters is paramount.
\end{abstract}

\begin{IEEEkeywords}
Google Earth Engine, Bi-LSTM, Topography, BIPAD, Multi-Scale Spatial Harmonization, 500m Resolution, Nepal, Multi-Hazard Prediction.
\end{IEEEkeywords}

\section{Introduction}
\IEEEPARstart{N}{epal} occupies a central position in the Himalayan arc, characterized by extreme topographic gradients and a monsoon-dominated climate. These factors contribute to a highly volatile environmental regime where hazards often occur in cascades. For example, extreme precipitation events frequently trigger landslides, which may subsequently block rivers and lead to catastrophic flash flooding (dam-burst floods).

Traditional disaster monitoring systems in the region have historically relied on sparse meteorological station networks and single-hazard retrospective mapping. However, the emergence of planetary-scale cloud computing platforms like Google Earth Engine (GEE) \cite{gorelick2017google} has enabled a transition toward proactive, data-driven multi-hazard prediction. 

This research proposes a unified multi-sensor fusion approach. By integrating active microwave sensors (Sentinel-1) for moisture and structural monitoring, multispectral optical sensors (Sentinel-2) for vegetation health, thermal sensors (MODIS) for land temperature, and atmospheric sensors (Sentinel-5P) for air quality, we aim to build a holistic "Environmental Risk Score." The core scientific objective is to evaluate the efficacy of the Long Short-Term Memory (LSTM) network \cite{hochreiter1997lstm}, a type of recurrent neural network, in predicting hazards that are driven by cumulative temporal factors.

\section{Literature Review}

\subsection{Multi-Hazard Susceptibility in the Himalayas}
The Himalayas represent one of the most active tectonic and geomorphological zones in the world. Petley et al. \cite{petley2007landslides} established that fatal landslides in Nepal are predominantly rainfall-induced, with a clear correlation between monsoon intensity and landslide frequency. However, single-factor threshold models often fail to account for the diversity of triggers, such as the unseasonal post-monsoon extreme rainfall events observed in October 2021 \cite{thapa2021melamchi}.

\subsection{Remote Sensing and GEE Applications}
The use of Sentinel-1 (SAR) has been transformative for flood monitoring due to its cloud-penetrating capability \cite{torres2012sentinel1}. Similarly, Sentinel-2 provides high-resolution data for computing the Normalized Difference Vegetation Index (NDVI), which is critical for assessing forest fire susceptibility and slope stability \cite{drusch2012sentinel}. Veefkind et al. \cite{veefkind2012sentinel5p} highlighted the potential of the TROPOMI sensor on Sentinel-5P for monitoring atmospheric composition, a feature we utilize for predicting air pollution and wildfire smoke impacts in the Kathmandu Valley and the Terai region.

\subsection{Deep Learning for Spatio-Temporal Prediction}
Machine learning models such as Random Forest (RF) and Support Vector Machines (SVM) are widely used for static susceptibility mapping. However, environmental hazards are fundamentally time-dependent. Wang et al. \cite{wang2019lstm} demonstrated that LSTM networks outperform traditional ML models in landslide prediction by maintaining an internal state that captures the history of antecedent precipitation and soil saturation. This research extends this methodology to a multi-hazard framework in Nepal.

\section{Study Area}
The study area encompasses the entire sovereign territory of Nepal, spanning from the lowland Terai plains ($60$~m above sea level) to the high Himalayan peaks ($>8,000$~m). The region is divided into five physiographic zones: Terai, Siwaliks, Middle Mountains, High Mountains, and High Himalaya. These zones exhibit diverse climatic patterns, from tropical to alpine, making Nepal a challenging yet ideal laboratory for testing multi-hazard prediction systems.

\section{Methodology}

\subsection{Data Acquisition and Preprocessing}
We utilized the GEE Python API to extract a unified daily time-series from 2021 to 2023. The data sources and their research applications are summarized in Table I.

\begin{table}[h!]
\centering
\caption{System Environmental Feature Matrix}
\begin{tabular}{l l l l}
\toprule
\textbf{Sensor} & \textbf{Dataset} & \textbf{Variable} & \textbf{Hazard Goal} \\
\midrule
Sentinel-1 & S1\_GRD & VV Polarization & Floods/Landslides \\
Sentinel-2 & S2\_SR\_HARMONIZED & B8, B4 (NDVI) & Fires/Landslides \\
Sentinel-5P & S5P NRTI L3 & Aerosol Index & Air Pollution \\
MODIS & MOD11A1 & LST\_Day\_1km & Drought/Fires \\
TerraClimate & TERRACLIMATE & Prcp, SoilM, Wind & General Trigger \\
\bottomrule
\end{tabular}
\end{table}

\subsection{Physics-Informed Feature Engineering}
To elevate the resolution and physical consistency, we implement two primary constraints:
\begin{enumerate}
    \item \textbf{Topographic Morphometry:} We extract static features from the 30m SRTM Digital Elevation Model, including Slope, Aspect, and the Topographic Position Index (TPI). These high-resolution features are spatially averaged over a 500m footprint using \texttt{ee.Reducer.mean()} at 30m computational scale, preserving the true statistical distribution of the terrain rather than point-sampling a single pixel.
    \item \textbf{Multi-Scale Spatial Harmonization:} Coarse-resolution climate variables (TerraClimate at $\sim$4km, MODIS LST at $\sim$1km, and Sentinel-5P at $\sim$7km) are bilinear-interpolated to the 500m analysis grid, while high-resolution data (Sentinel-1 at 10m, Sentinel-2 at 10m, SRTM at 30m) are spatially aggregated via mean reduction. This dual strategy prevents aliasing artifacts from high-resolution sources while smoothly upsampling low-resolution climate fields.
\end{enumerate}

\subsection{Ground Truth Integration (BIPAD)}
The model transitions from synthetic proxies to official disaster records. We utilize the Nepal BIPAD portal database, extracting 873 documented 'Landslide' and 'Flood' events (2021--2023). Each event is mapped to a 500m radius on the date of occurrence. For training, a balanced sampling approach (1:5 ratio) is employed, comparing hazard locations against stable regions across the country.

\subsection{Bi-Directional Deep Learning Architecture}

\subsubsection{Temporal Window Justification}
The system constructs 7-day multi-channel sequences $[F_{T-6}, \dots, F_T]$ for each sample location. The 7-day window was selected as a principled compromise between the competing temporal scales of the target hazards: flash floods depend on short-term precipitation intensity (1--3 days), while rainfall-induced landslides correlate with 10--30 day moisture accumulation \cite{petley2007landslides}. A critical design consideration is that TerraClimate soil moisture and precipitation variables inherently encode longer-term accumulation---the extraction pipeline uses a 45-day lookback window when querying TerraClimate composites, meaning each daily feature vector already carries an implicit signal of antecedent saturation extending well beyond the 7-day sequence boundary. The 7-day LSTM window thus operates over features that implicitly represent $\sim$45 days of hydrological memory.

\subsubsection{Why Bi-LSTM?}
The Bi-directional LSTM processes the 7-day \textit{historical} window in both forward and backward directions. Crucially, this does \textbf{not} require future data beyond $T_0$---at inference time, the model receives only the past 7 days $[T_{-6}, \dots, T_0]$ and processes this fixed window bidirectionally. The backward pass allows the network to contextualize early-window features (e.g., a rainfall spike at $T_{-6}$) in light of the current state at $T_0$ (e.g., saturated soil), enabling it to identify \textit{which} antecedent conditions were causally relevant. A standard unidirectional LSTM processes $T_{-6}$ without knowledge of the eventual outcome at $T_0$, limiting its ability to learn such retrospective causal patterns.

\subsubsection{Environmental Hazard Risk Score}
The final Dense layer with sigmoid activation produces the \textbf{Environmental Hazard Risk Score} $\hat{P}(H | \mathbf{X})$, defined as:
\begin{equation}
    \hat{P}(H | \mathbf{X}) = \sigma\left(\mathbf{W}^\top \cdot h_{\text{BiLSTM}} + b\right) \in [0, 1]
\end{equation}
where $h_{\text{BiLSTM}}$ is the concatenated hidden state from the forward and backward LSTM passes, $\sigma$ is the sigmoid function, and $\mathbf{X} = [F_{T-6}, \dots, F_T]$ is the 7-day input sequence. A score $> 0.5$ classifies the location--date pair as a hazard event; the continuous score itself serves as a normalized risk index suitable for early warning systems.

The model uses a weighted Binary Cross Entropy loss to address the extreme class imbalance:
\begin{equation}
    \mathcal{L} = -\frac{1}{N}\sum_{i=1}^{N}\left[w_1 \cdot y_i \log(\hat{p}_i) + w_0 \cdot (1-y_i)\log(1-\hat{p}_i)\right]
\end{equation}
where $w_1 = 3.08$ and $w_0 = 0.60$ are inverse-frequency class weights computed from the training set.

\section{Results}

\subsection{Model Comparison and Performance}
To justify the architectural choice of Bi-LSTM over simpler alternatives, we evaluate four models on the \textit{identical} train/test split (80/20, $n=5{,}188$). Tree-based baselines (Random Forest, Gradient Boosting) operate on the flattened 70-dimensional feature vector ($7 \times 10$), while recurrent models (LSTM, Bi-LSTM) process the native 3D temporal structure.

\begin{table}[h!]
\centering
\caption{Comparative Model Performance on Nepal Hazard Dataset}
\begin{tabular}{l c c c c c}
\toprule
\textbf{Model} & \textbf{Acc.} & \textbf{Prec.} & \textbf{Recall} & \textbf{F1} & \textbf{AUC} \\
\midrule
Random Forest       & 0.90 & 0.76 & 0.70 & 0.73 & 0.948 \\
Gradient Boosting   & 0.91 & 0.81 & 0.71 & 0.76 & 0.953 \\
LSTM (Unidirectional) & 0.87 & 0.63 & 0.85 & 0.72 & 0.930 \\
\textbf{Bi-LSTM}    & \textbf{0.88} & \textbf{0.66} & \textbf{0.81} & \textbf{0.73} & \textbf{0.933} \\
\bottomrule
\end{tabular}
\end{table}

\subsection{Analysis of Trade-offs}

The results reveal a nuanced trade-off between precision-dominant and recall-dominant architectures:

\begin{itemize}
    \item \textbf{Tree-based models} (RF, GB) achieve higher overall accuracy (90--91\%) and precision (76--81\%) but suffer from lower recall (70--71\%). In a disaster early-warning context, this means they \textit{miss} $\sim$30\% of real hazard events---an unacceptable false negative rate for life-safety applications.
    \item \textbf{Recurrent models} (LSTM, Bi-LSTM) sacrifice $\sim$3\% overall accuracy but achieve substantially higher recall (81--85\%). The Bi-LSTM, compared to the unidirectional variant, improves precision from 0.63 to 0.66 (+4.8\%) without significant recall loss, demonstrating the value of bidirectional context within the historical window.
    \item \textbf{Why not tune Gradient Boosting for recall?} One could lower GB's classification threshold to increase recall. However, this trades precision catastrophically: because tree-based models do not inherently model temporal dependencies and instead rely on engineered lag features, they lack the structural inductive bias to distinguish \textit{when} a feature pattern is dangerous. The Bi-LSTM's high recall emerges naturally from its sequential architecture---it learns that ``5 days of moderate rain on a 25$^\circ$ slope followed by a precipitation spike'' is qualitatively different from the same cumulative total distributed uniformly. GB, operating on a flattened 70-dimensional vector, must discover such interactions combinatorially rather than structurally.
\end{itemize}

For an early warning system where \textbf{missing a real disaster is far worse than a false alarm}, the Bi-LSTM's recall-dominant profile (0.81 recall, 0.66 precision) is operationally preferable to Gradient Boosting's precision-dominant profile (0.71 recall, 0.81 precision).

\subsection{Validation: BIPAD Incident Correlation}
By validating against official 2021--2023 BIPAD records, we established that the model successfully differentiates between ``Stable'' high-moisture periods and actual ``Incident Trigger'' events. The use of a weighted loss function allowed the architecture to overcome the extreme class imbalance of disaster datasets, maintaining high recall even when the frequency of ``Landslide'' or ``Flood'' events is low.

\subsection{Error Analysis: Characterizing Missed Events}
To evaluate the operational safety of the system, we perform a deep characterization of the 37 missed hazard events (False Negatives). By cross-referencing these incidents with BIPAD severity metadata, we find that the model overwhelmingly fails on minor, low-impact events rather than major disasters.

\begin{itemize}
    \item \textbf{Casualty Distribution:} Of the 37 missed events, 34 (92\%) resulted in zero deaths or injuries, representing property-only or minor incidents. Only three events involved casualties (1 death, 2 missing, 2 injured in total), compared to 162 successfully detected hazard events.
    \item \textbf{Precipitation Bias:} Feature analysis reveals that missed events have a mean precipitation of 162.4\,mm at $T_0$, which is 53.8\% lower than the 351.6\,mm mean for detected events. This suggests the model has learned a robust ``wet-trigger'' threshold, but remains blind to dry-season failures.
    \item \textbf{Case Analysis:} The single fatal missed event (Dhading, Jan 2022) occurred during the dry winter season (13\,mm rain), likely triggered by progressive geological weathering or seismic activity rather than monsoon precipitation. Another missed casualty event occurred on a 2.5$^\circ$ slope, falling outside the model's learned geomorphological profile for steep-slope failures.
\end{itemize}

These findings indicate that while the Bi-LSTM is a reliable monsoon-season predictor, its primary failure mode is the detection of dry-season geologically-driven hazards which lack a discernible satellite-based meteorological signature.

\section{Discussion: The ``Physical-AI'' Synergy}
The research underscores the necessity of ``Physically-Constrained'' architectures. Deep learning alone, without topographic context, often fails to distinguish between heavy rainfall on a stable plain versus a vulnerable 30-degree slope. By ``injecting'' the 30m SRTM physics into the 500m temporal sequence, we achieved a more robust and geomorphologically sound prediction engine.

The comparative analysis (Table II) demonstrates that no single model dominates across all metrics. Tree-based models excel at precision through learned feature interactions but do not inherently model temporal dependencies. Recurrent models capture sequential moisture buildup through their structural inductive bias but require richer feature engineering to match tree-based precision. The Bi-LSTM represents the best compromise for operational early warning, where recall is the primary objective.

\subsection{Spatial Validation}
Fig.~1 presents a spatial overlay of the Bi-LSTM's predicted hazard risk scores against BIPAD ground truth events across Nepal's test set ($n=1{,}038$). The map reveals that high-probability predictions (warm colors) cluster in known hazard-prone corridors---the Middle Hills and Siwalik foothills---while stable Terai and high Himalayan regions receive consistently low risk scores. Critically, the False Negative markers (missed events) do not concentrate in any single geographic zone, suggesting the model generalizes across Nepal's physiographic diversity rather than over-fitting to a specific region.

\begin{figure}[h!]
    \centering
    \includegraphics[width=\columnwidth]{spatial_risk_map.png}
    \caption{Spatial distribution of Bi-LSTM hazard risk predictions vs. BIPAD ground truth on the test set. Red circles = correctly detected hazards (True Positives); cyan X-markers = missed events (False Negatives); orange squares = false alarms.}
    \label{fig:spatial_map}
\end{figure}

\section{Limitations and Future Work}
\begin{enumerate}
    \item \textbf{Temporal Window Extension:} The current 7-day sequential window may be insufficient for slow-onset landslides driven by 10--30 day rainfall accumulation. While TerraClimate's implicit 45-day lookback partially compensates, future work should extend the explicit LSTM window to 14--30 days and evaluate impact on recall for deep-seated landslide events.
    \item \textbf{Hybrid Architecture:} A hybrid approach combining Gradient Boosting's precision with Bi-LSTM's recall (e.g., stacking or attention-based fusion) may yield superior performance across both metrics.
    \item \textbf{Higher-Resolution DEM:} Integration of sub-10m LiDAR or ALOS PALSAR DEM data could improve localization of slope instability beyond the current 30m SRTM baseline.
    \item \textbf{Spatial Cross-Validation:} The current random train/test split does not account for spatial autocorrelation. Future iterations should implement blocked spatial cross-validation to ensure geographic generalization.
\end{enumerate}

\section{Conclusion}
This research presents a physically-constrained, multi-sensor, temporally-aware hazard prediction system optimized for recall in disaster scenarios. By fusing six satellite data streams with 30m SRTM topographic physics into a Bi-LSTM architecture, we produce a continuous Environmental Hazard Risk Score $\hat{P}(H | \mathbf{X}) \in [0,1]$ for any 500m location--date pair across Nepal. Comparative evaluation across four model architectures confirms that the Bi-LSTM achieves the highest hazard recall (0.81) while maintaining competitive AUC (0.933), making it operationally suitable for early warning systems where minimizing missed disasters is paramount. The spatial validation (Fig.~1) demonstrates geographic generalization across Nepal's diverse physiographic zones, and the system's modular GEE-based extraction pipeline makes it directly transferable to other Himalayan nations facing similar multi-hazard challenges.

\bibliographystyle{IEEEtran}
\bibliography{references}

\end{document}

